% ---------------------------------------------------------------------------
% appendix-quickref.tex -- design-side quick reference.
% Split off from the original SystemVerilog cookbook's appendix-cheatsheet.tex
% when the book was divided into companion booklets; the command-line
% reference moved to the Simulation Cookbook (see backmatter/appendix-commands.tex
% there).
% ---------------------------------------------------------------------------
\chapter{Quick Reference}
\label{app:designref}

\chapabstract{The VHDL/SystemVerilog language mapping, compressed onto a few
pages to keep next to the keyboard. Everything here is explained properly
somewhere in the preceding chapters; this is the version you use once you
already understand it.}

\section{The ten things that will bite you first}

\begin{enumerate}
  \item \textbf{Widths do not have to match.} \code{a + b} with different
        widths compiles. The result is truncated or extended silently. Lint
        for it.
  \item \textbf{One unsigned operand makes the whole expression unsigned.}
        Mixing \code{logic signed} and \code{logic} in one expression
        produces an unsigned result, and your negative numbers become very
        large positive ones.
  \item \textbf{\code{>>} is not \code{>>>}.} On a signed value you almost
        always want \code{>>>}. VHDL's \code{shift\_right} on a
        \code{signed} is arithmetic by definition and never wrong.
  \item \textbf{\code{<=} in an \code{always\_comb} block, or \code{=} in an
        \code{always\_ff} block}, will simulate and may even synthesise. It
        is still wrong. Non-blocking for sequential, blocking for
        combinational.
  \item \textbf{A function is static by default.} Declare
        \code{function automatic} unless you know why you want otherwise.
        The same applies to tasks and to variables inside forked blocks.
  \item \textbf{An undeclared identifier becomes a one-bit wire.} Put
        \code{`default\_nettype none} at the top of every file and
        \code{`default\_nettype wire} at the bottom.
  \item \textbf{\code{`timescale} is positional and inherited.} A module
        with no preceding directive picks up whatever the last compiled file
        set. Put one in every file, or set it on the command line.
  \item \textbf{A class handle is a reference.} \code{b = a} makes two names
        for one object. Use \code{copy()} or a \code{new} plus field copy.
  \item \textbf{\code{randomize()} returns a status you must check.} An
        unsatisfiable constraint set silently leaves the old values in place
        if you ignore it.
  \item \textbf{An unsized literal is 32 bits.} \code{'1} is all ones at the
        context width; \code{1} is a 32-bit signed one. They are not the
        same thing in a concatenation.
\end{enumerate}

\section{Construct map}

\begin{longtable}{@{}p{0.32\textwidth}p{0.34\textwidth}p{0.28\textwidth}@{}}
\toprule
\textbf{VHDL} & \textbf{SystemVerilog} & \textbf{Note}\\
\midrule
\endfirsthead
\toprule
\textbf{VHDL} & \textbf{SystemVerilog} & \textbf{Note}\\
\midrule
\endhead
\code{entity} + \code{architecture} & \code{module} \ldots\ \code{endmodule}
  & one unit, not two\\
\code{library}/\code{use} & \code{import pkg::*} & plus \code{`include}\\
\code{generic} & \code{parameter} & \code{localparam} if not overridable\\
\code{std\_ulogic} & \code{logic} & 4-state, unresolved\\
\code{std\_logic} & \code{wire} & resolved; use \code{logic} unless there are two drivers\\
\code{std\_logic\_vector(7 downto 0)} & \code{logic [7:0]} & \\
\code{unsigned}/\code{signed} & \code{logic [N-1:0]} /
  \code{logic signed [N-1:0]} & no separate type\\
\code{(others => '0')} & \code{'0} & context-width fill\\
\code{a \& b} & \code{\{a, b\}} & concatenation\\
--- & \code{\{N\{a\}\}} & replication\\
\code{process (all)} & \code{always\_comb} & \code{always\_comb} also runs at
  time~0\\
\code{process (clk)} + \code{rising\_edge} & \code{always\_ff @(posedge clk)}
  & \\
\code{variable} in a process & blocking \code{=} & lifetime differs\\
\code{signal} assignment & non-blocking \code{<=} & \\
\code{wait for 10 ns} & \code{\#10ns} & bare \code{\#10} takes its unit from \code{`timescale}\\
\code{wait until cond} & \code{wait(cond)} or \code{@(...)} & \\
\code{for \ldots generate} & \code{for (genvar \ldots)} & \code{generate}
  optional\\
\code{if \ldots generate} & \code{if (\ldots) begin : name} & name the block\\
\code{port map (a => x)} & \code{.a(x)} & \code{.*} exists, use sparingly\\
\code{function} & \code{function automatic} & \\
\code{procedure} & \code{task} or \code{void function} & \\
\code{to\_integer(unsigned(v))} & \code{int'(v)} & implicit in most contexts\\
\code{resize(v, N)} & assignment or cast & implicit and silent\\
\code{shift\_right(s, n)} & \code{s >>> n} & signed operand required\\
\code{'length}, \code{'range} & \code{\$bits}, \code{\$size} & no
  \code{'range} equivalent\\
\code{'event}, \code{'stable} & \code{\$rose}, \code{\$fell},
  \code{\$stable} & sampled semantics\\
\code{assert \ldots severity error} & \code{assert (\ldots) else \$error(...)}
  & \\
\code{report} & \code{\$display} & \\
\code{integer'image(x)} & \code{\%0d} in a format string & \\
\code{std.textio} & \code{\$fopen}/\code{\$fscanf}/\code{\$fdisplay} & \\
\code{record} & \code{struct packed} & packed struct is a bit vector\\
\code{type \ldots is (A, B)} & \code{typedef enum logic [1:0] \{A, B\}} &
  encoding is explicit\\
\code{protected type} & \code{class} & class is not synthesisable\\
\code{configuration} & \code{bind} / factory override & no direct equivalent\\
--- & \code{interface} / \code{modport} & no VHDL-2008 equivalent\\
--- & \code{covergroup}, \code{constraint}, \code{class} & verification only\\
\code{ieee.fixed\_pkg} & --- & track the binary point yourself\\
\bottomrule
\end{longtable}

\section{Format string cheat sheet}

\begin{center}
\small
\begin{tabular}{@{}llll@{}}
\toprule
\textbf{Specifier} & \textbf{Meaning} & \textbf{Specifier} & \textbf{Meaning}\\
\midrule
\code{\%d} & decimal, padded      & \code{\%0d} & decimal, minimum width\\
\code{\%h} & hexadecimal          & \code{\%0h} & hexadecimal, no padding\\
\code{\%b} & binary               & \code{\%o}  & octal\\
\code{\%s} & string               & \code{\%c}  & character\\
\code{\%f} & real, fixed          & \code{\%e}  & real, exponential\\
\code{\%g} & real, shorter of the two & \code{\%t} & time, per
  \code{\$timeformat}\\
\code{\%p} & any aggregate, pretty & \code{\%m} & hierarchical name\\
\code{\%v} & net strength         & \code{\%\%} & a literal per cent sign\\
\bottomrule
\end{tabular}
\end{center}

\section{Severity and what it does}

\begin{center}
\small
\begin{tabular}{@{}lll@{}}
\toprule
\textbf{VHDL} & \textbf{SystemVerilog} & \textbf{Effect}\\
\midrule
\code{severity note}    & \code{\$info}    & message only\\
\code{severity warning} & \code{\$warning} & message only\\
\code{severity error}   & \code{\$error}   & message, counts as a failure\\
\code{severity failure} & \code{\$fatal(n)} & message, ends the simulation
  with status \code{n}\\
\bottomrule
\end{tabular}
\end{center}
